
\chapter{Principal Stratification}
 \label{chapter::principal-stratification}

本书的 \cref{part::rcts}--\cref{part::instrumentalvariables} 关注 treatment 对 outcome 的因果效应，并且有时会调整一些 treatment 之前观测到的 covariates。许多应用中还会出现一个 post-treatment variable $M$，它发生在 treatment 之后，并且与 outcome 有关。一个重要问题是：如何恰当地使用这个 post-treatment variable $M$。本章先从几个动机例子开始，然后介绍 \citet{frangakis2002principal} 基于 potential outcomes 对这个问题给出的 principal stratification 表述。
 


\section{Motivating Examples}



\begin{example}[noncompliance]\label{eg::noncompliance}
在存在 noncompliance 的随机实验中，我们可以用 $M$ 表示实际接受的 treatment；它受 treatment assignment $Z$ 影响，并且影响 outcome $Y$。
在这个例子中，$M$ 与 \cref{chapter::iv-experiment} 中的 $D$ 含义相同。
\end{example}


\begin{example}[truncation by death]\label{eq::truncationbydeath}
在针对重症患者的随机实验中，一些患者可能在 outcome $Y$ 被测量之前死亡，例如在测量生活质量之前死亡。这个例子中的 post-treatment variable $M$ 是 survival status 的二元指示变量。
\end{example}




\begin{example}[unemployment]\label{eg::unemployment}
在就业培训项目中，单位被随机分配到 treatment group 与 control group，并报告其 employment status $M$ 和 wage $Y$。此时 post-treatment variable 是 employment status $M$ 的二元指示变量。
\end{example}




\begin{example}
[surrogate endpoint]\label{eg::surrogate-endpoints}
在临床试验中，真正关心的 outcomes（例如 30 年 survival）需要长期而昂贵的随访。实践者因此会在随访早期收集一些更容易测量的变量。这些变量称为 ``surrogate endpoints''。一个具体例子来自 HIV 患者的临床试验，其中候选 surrogate endpoint $M$ 是 CD4 cell count（回忆一下，CD4 cells 是抵抗感染的 white blood cells）。
\end{example}



上面的 \cref{eg::noncompliance}--\cref{eg::surrogate-endpoints} 有一个共同点：变量 $M$ 发生在 treatment 之后，并且与 outcome 相关。$M$ 可能位于从 $Z$ 到 $Y$ 的 causal pathway 上；\cref{fig::intermediate-diagram}(a) 描绘了这种机制。\cref{eg::noncompliance} 对应 \cref{fig::intermediate-diagram}(a)。$M$ 也可能不在从 $Z$ 到 $Y$ 的 causal pathway 上；\cref{fig::intermediate-diagram}(b) 描绘了这种机制。\cref{eq::truncationbydeath} 和 \cref{eg::unemployment} 对应 \cref{fig::intermediate-diagram}(b)。\cref{eg::surrogate-endpoints} 可能对应 \cref{fig::intermediate-diagram}(a)，也可能对应 (b)，这取决于 surrogate endpoint 的选择。

 


\begin{figure}[h]
\centering
\begin{subfigure}[]{\textwidth}
                \centering
$$
\begin{xy}
\xymatrix{
&& & U \ar[dl]   \ar[dr] \\
Z \ar[rr] \ar@/_1.5pc/[rrrr] & & M\ar[rr] & & Y\\
}
\end{xy}
$$ 
               \caption{$M$ 位于从 $Z$ 到 $Y$ 的 causal pathway 上，其中 $Z$ 是 randomized，$U$ 表示 unmeasured confounding}
        \end{subfigure}%
        
        \begin{subfigure}[]{\textwidth}
                \centering
$$
\begin{xy}
\xymatrix{
&& & U \ar[dl]   \ar[dr] \\
Z \ar[rr] \ar@/_1.5pc/[rrrr] & & M  & & Y\\
}
\end{xy}
$$
                \caption{$M$ 不位于从 $Z$ 到 $Y$ 的 causal pathway 上，其中 $Z$ 是 randomized，$U$ 表示 unmeasured confounding}
        \end{subfigure}%
        
\caption{带有 post-treatment variable $M$ 的 causal diagrams}\label{fig::intermediate-diagram}
\end{figure}


在实践中，底层 causal diagrams 可能比 \cref{fig::intermediate-diagram} 中的图复杂得多。本章采用 \citet{frangakis2002principal} 的表述；这一表述不假设具体的底层 causal diagrams。





\section{The Problem of Conditioning on the Post-Treatment Variable}

处理 post-treatment variable $M$ 的一个 naive 方法，是像处理 pretreatment covariate 那样，对其观测值进行条件化。然而，$M$ 与 $X$ 有根本差异，因为一般来说 $M$ 会受到 treatment 影响，而 $X$ 不会。数据分析者在评估 treatment 对 outcome 的 average causal effect 时，不应 conditioning on any post-treatment variables，这也是一个常见的 ``rule of thumb'' \citep{cochran1957analysis, rosenbaum1984consquences}。基于 potential outcomes，\citet{frangakis2002principal} 给出了如下深刻解释。

为简单起见，本章聚焦于 CRE。

\begin{assumption}
[CRE with an intermediate variable]
\label{assume::complete-randomization-withM}
我们有
$$
Z \ind  \{ M(1), M(0), Y(1), Y(0), X \} . 
$$
\end{assumption}

 
Conditioning on $M= m$，我们比较
\begin{equation}
\pr(Y\mid Z=1, M=m)\label{eq::proby|z1m}
\end{equation}
和
\begin{equation}
\pr(Y\mid Z=0, M=m).\label{eq::proby|z0m}
\end{equation}
这个比较看起来很直观，因为它在 post-treatment variable 取同一值时，度量 treated group 与 control group 中 outcome distributions 的差异。当 $M$ 是 pre-treatment covariate 时，这个比较会给出合理的 subgroup effect。然而，当 $M$ 是 post-treatment variable 时，这个比较的解释就有问题。在 \cref{assume::complete-randomization-withM} 下，我们可以将 \cref{eq::proby|z1m,eq::proby|z0m} 中的概率改写为
\begin{eqnarray*}
\pr(Y\mid Z=1, M=m)  &=& \pr\{ Y(1)\mid Z=1, M(1) = m  \} \\
&=& \pr\{ Y(1)\mid   M(1) = m  \}
\end{eqnarray*}
以及
\begin{eqnarray*}
\pr(Y\mid Z=0, M=m) &=&  \pr\{ Y(0)\mid Z=0, M(0) = m  \} \\
&=& \pr\{ Y(0)\mid   M(0) = m  \} . 
\end{eqnarray*}
因此，在 CRE 下，比较 \cref{eq::proby|z1m,eq::proby|z0m} 等价于比较不同单位子集上的 $Y(1)$ 与 $Y(0)$ 的分布；原因是如果 $Z$ 影响 $M$，那么满足 $M(1)=m$ 的单位一般不同于满足 $M(0)=m$ 的单位。于是，除非 $M(1)=M(0)$，否则 conditioning on $M=m$ 的比较通常没有 causal interpretation。\footnote{从 causal diagrams 出发也能得到同样结论。
在 \cref{fig::intermediate-diagram} 中，尽管由 $Z$ 的 randomization 得到 $Z\ind U$，conditioning on $M$ 会引入 ``collider bias''，从而导致 $Z\nind U$。
} 


重新看 \cref{eg::noncompliance}。在 monotonicity assumption $M(1)\geq M(0)$ 下，比较 $\pr(Y\mid Z=1, M=1)$ 与 $\pr(Y\mid Z=0, M=1)$，等价于将 compliers 与 always-takers 的 treated potential outcomes，同 always-takers 的 control potential outcomes 作比较。\cref{problem::as-treated-analysis} 的第 3 部分已经指出了这种分析的缺陷。


重新看 \cref{eq::truncationbydeath}。如果 treatment 改善 survival status，那么 treatment 可能比 control 挽救更多身体较弱的患者。在这种情形下，满足 $M(1)=1$ 的单位比满足 $M(0)=1$ 的单位更弱，因此 naive comparison 会产生偏向 control 的有偏结果。

\section{Conditioning on the Potential Values of the Post-Treatment Variable}


\citet{frangakis2002principal} 提出 conditioning on post-treatment variable 的 joint potential value $U = \{ M(1), M(0) \}$，并比较
$$
\pr\{  Y(1) \mid M(1) = m_1, M(0) = m_0 \}
$$
和
$$
\pr\{  Y(0) \mid M(1) = m_1, M(0) = m_0 \}
$$
其中 $(m_1,m_0)$ 固定。这是在同一个单位子集内比较 treatment 与 control 下的 potential outcomes，该子集满足 $M(1)=m_1$ 且 $M(0)=m_0$。\citet{frangakis2002principal} 将这种策略称为 {\it principal stratification}，并把 $\{M(1),M(0)\}$ 看作一个 pre-treatment covariate。基于这个思想，我们可以定义
$$
\tau(m_1, m_0) = E\{  Y(1) - Y(0) \mid M(1) = m_1, M(0) = m_0\}
$$
作为满足 $M(1)=m_1, M(0)=m_0$ 的 subgroup 的 principal stratification average causal effect。对于二元 $M$，我们有四个 subgroups：
\begin{equation}
\label{eq::psace}
\left\{ \begin{array}{rcl}
\tau(1,1) &=& E\{  Y(1) - Y(0) \mid M(1) = 1, M(0) =1\} ,\\
\tau(1,0) &=&E\{  Y(1) - Y(0) \mid M(1) = 1, M(0) =0\},\\
\tau(0,1) &=&E\{  Y(1) - Y(0) \mid M(1) = 0, M(0) =1\},\\
\tau(0,0) &=&E\{  Y(1) - Y(0) \mid M(1) = 0, M(0) = 0 \}.
\end{array} 
\right.
\end{equation}
由于 $\{M(1),M(0)\}$ 不受 treatment 影响，它是一个 covariate，因此 $\tau(m_1,m_0)$ 是 subgroup causal effect。对于满足 $M(1)=M(0)$ 的 subgroups，treatment 不改变 intermediate variable，所以 $\tau(1,1)$ 与 $\tau(0,0)$ 度量 {\it dissociative effects}。
对于其他满足 $m_1\neq m_0$ 的 subgroups，principal stratification average causal effects $\tau(m_1,m_0)$ 度量 {\it associative effects}。这些术语来自 \citet{frangakis2002principal}；它并不假设 $M$ 位于从 $Z$ 到 $Y$ 的 causal pathway 上。当图形是 \cref{fig::intermediate-diagram}(a) 时，我们可以将 dissociative effects 解释为 $Z$ 对 $Y$ 的 direct effects，并且这些 effects 独立于 $M$ 起作用；但我们不能简单地将 associative effects 解释为 $Z$ 对 $Y$ 的 direct 或 indirect effects。

 

\setcounter{example}{0}
\begin{example}[noncompliance]
在 noncompliance 下，\cref{eq::psace} 包含 always takers、compliers、defiers 和 never takers 的 average causal effects \citep{imbens1994identification, angrist1996identification}。
\end{example}



\begin{example}[truncation by death]
由于 outcome 只有在患者存活时才有良好定义，\cref{eq::psace} 中三个 subgroup causal effects 没有意义，唯一良好定义的 subgroup effect 是
\begin{eqnarray}
\label{eq::survivor}
\tau(1,1) = E\{  Y(1) - Y(0) \mid M(1) = 1, M(0) =1\} .
\end{eqnarray}
它称为 survivor average causal effect \citep{rubin2006causal}。它是在无论 treatment status 如何都会存活的单位中，treatment 对 outcome 的 average causal effect。
\end{example}



\begin{example}[unemployment]
unemployment 问题与 truncation-by-death 问题同构，因为 wage 只有在单位首先被雇佣时才有良好定义。因此，唯一良好定义的 subgroup effect 是 \cref{eq::survivor}，即 employed average causal effect。此前，\citet{heckman1979sample} 提出了一个模型来处理工资建模中的 unemployment；这个模型现在称为 Heckman Selection Model，它把 unemployed 人群的 wages 看作 missing values\footnote{Heckman 在 2000 年获得诺贝尔经济学奖，获奖理由是 ``for his development of theory and methods for analyzing selective samples.''
他的模型包含两个阶段。第一，employment status 由一个 latent linear model 决定：
$$
M_i = 1(X_i\tran \beta + u_i \geq 0).
$$
第二，latent log wage 由一个线性模型决定：
$$
Y_i^* = W_i\tran \gamma + v_i
$$
并且只有当 $M_i=1$ 时，$Y_i^*$ 才作为 $Y_i$ 被观测到。在他的 two-stage model 中，covariates $X_i$ 与 $W_i$ 可以不同，errors $(u_i,v_i)$ 服从相关的 bivariate Normal。
}。然而，\citet{zhang2003estimation} 和 \citet{zhang2009likelihood} 认为，在 potential outcomes framework 下，$\tau(1,1)$ 是更有意义的量。
\end{example} 



\begin{example}[surrogate endpoint]
直观地说，我们希望通过 treatment 对 surrogate endpoint 的 effect 来评估 treatment 对 outcome 的 effect。因此，一个好的 surrogate endpoint 应该满足两个条件：第一，如果 treatment 不影响 surrogate，那么它也不影响 outcome；第二，如果 treatment 影响 surrogate，那么它也影响 outcome。第一个条件被 \citet{frangakis2002principal} 称为 ``causal necessity''，第二个条件被 \citet{gilbert2008evaluating} 称为 ``causal sufficiency''。基于二元 surrogate endpoint 的 \cref{eq::psace}，causal necessity 要求 $\tau(1,1)$ 和 $\tau(0,0)$ 为零，而 causal sufficiency 要求 $\tau(1,0)$ 和 $\tau(0,1)$ 非零。
\end{example}





\section{Statistical Inference and Its Difficulty}
\label{sec::difficulty-of-PS}

在 \cref{eg::noncompliance} 中，如果有 randomization、monotonicity 和 exclusion restriction，那么我们可以识别 complier average causal effect $\tau(1,0)$。这是 \cref{chapter::iv-experiment} 推导出的关键结果。


然而，在其他例子中，我们不能施加 exclusion restriction assumption。例如，$\tau(1,1)$ 是 \cref{eq::truncationbydeath,eg::unemployment} 中的主要目标参数，而 $\tau(1,1)$ 和 $\tau(0,0)$ 都是 \cref{eg::surrogate-endpoints} 中感兴趣的参数。没有 exclusion restriction assumption 时，识别 principal stratification average causal effect 非常困难。有时，我们甚至不能施加 monotonicity assumption，因此一开始就无法识别 latent strata 的比例。




\subsection{Special case: truncation by death with binary outcome}
我用一个具有 binary treatment、binary survival status 和 binary outcome 的简单设定来说明 principal stratification 下 statistical inference 的思想，尤其是其困难。


除了 \cref{assume::complete-randomization-withM} 之外，我们还施加 monotonicity。

\begin{assumption}
[monotonicity]\label{assume::monotonicity-ps}
$M(1) \geq M(0)$.
\end{assumption}


\cref{thm::cace-mixture-components} 表明，在 \cref{assume::complete-randomization-withM,assume::monotonicity-ps} 下，我们可以识别三个 latent strata 的比例：
\begin{eqnarray*} 
 \pi_{(1,1)} &=& \pr(M=1\mid Z=0), \\
 \pi_{(0,0)} &=& \pr(M=0\mid Z=1), \\
 \pi_{(1,0)} &=& \pr(M=1\mid Z=1) - \pr(M=1\mid Z=0). 
\end{eqnarray*} 


我们的目标是识别 survivor average causal effect $\tau(1,1)$。首先，我们可以很容易识别 $E\{Y(0)\mid M(1)=1,M(0)=1\}$，因为观测组 $(Z=0,M=1)$ 只由 survivors 构成：
$$
E\{ Y(0)\mid M(1) = 1, M(0)=1 \} = E(Y\mid Z=0, M=1). 
$$
关键在于识别 $E\{Y(1)\mid M(1)=1,M(0)=1\}$。观测组 $(Z=1,M=1)$ 是两个 strata $(1,1)$ 和 $(1,0)$ 的 mixture，因此我们有
\begin{eqnarray} 
E(Y\mid Z=1,M=1) 
&=& \frac{ \pi_{(1,1)}}{ \pi_{(1,1)} + \pi_{(1,0)}} E\{  Y(1)\mid   M(1)=1, M(0)=1 \} \nonumber  \\
&&+ \frac{ \pi_{(1,0)}}{ \pi_{(1,1)} + \pi_{(1,0)}} E\{  Y(1)\mid  M(1)=1, M(0)=0 \}. \nonumber  \\
&& \label{eq::bound-equation-basic}
\end{eqnarray} 
\cref{eq::bound-equation-basic} 中有两个未知参数，却只有一个方程。因此，我们不能由 \cref{eq::bound-equation-basic} 唯一确定 $E\{Y(1)\mid M(1)=1,M(0)=1\}$。不过，\cref{eq::bound-equation-basic} 仍然包含关于目标量的信息。
也就是说，按照 \cref{def::partial-identification}，$E\{Y(1)\mid M(1)=1,M(0)=1\}$ 是 partially identified。



对于二元 outcome $Y$，我们知道 $E\{Y(1)\mid M(1)=1,M(0)=0\}$ 被限制在 0 和 1 之间。因此，$E\{Y(1)\mid M(1)=1,M(0)=1\}$ 被限制在下面两个方程的解之间：
\begin{eqnarray*}
E(Y\mid Z=1,M=1) &=& \frac{ \pi_{(1,1)}}{ \pi_{(1,1)} + \pi_{(1,0)}} E\{  Y(1)\mid   M(1)=1, M(0)=1 \}  \\
&&\qquad + \frac{ \pi_{(1,0)}}{ \pi_{(1,1)} + \pi_{(1,0)}}  
\end{eqnarray*}
以及
$$
E(Y\mid Z=1,M=1) = \frac{ \pi_{(1,1)}}{ \pi_{(1,1)} + \pi_{(1,0)}} E\{  Y(1)\mid   M(1)=1, M(0)=1 \} .
$$
因此，$E\{Y(1)\mid M(1)=1,M(0)=1\}$ 的 lower bound 为
$$
\frac{  \{  \pi_{(1,1)} + \pi_{(1,0)} \} E(Y\mid Z=1,M=1) - \pi_{(1,0)}}{ \pi_{(1,1)}},
$$
upper bound 为
$$
\frac{  \{  \pi_{(1,1)} + \pi_{(1,0)} \} E(Y\mid Z=1,M=1)  }{ \pi_{(1,1)}} . 
$$
于是可以推出 $\tau(1,1)$ 的 bounds，总结在下面的 \cref{thm::bounds-on-SACE-binary-Y} 中。



\begin{theorem}
\label{thm::bounds-on-SACE-binary-Y}
在 \cref{assume::complete-randomization-withM,assume::monotonicity-ps} 下，若 $Y$ 为二元变量，则有
\begin{eqnarray*}
&&\frac{  \{  \pi_{(1,1)} + \pi_{(1,0)} \} E(Y\mid Z=1,M=1) - \pi_{(1,0)}}{ \pi_{(1,1)}} - E(Y\mid Z=0, M=1) \\
&\leq & \tau(1,1) \\
& \leq &\frac{  \{  \pi_{(1,1)} + \pi_{(1,0)} \} E(Y\mid Z=1,M=1)  }{ \pi_{(1,1)}}  - E(Y\mid Z=0, M=1). 
\end{eqnarray*}
\end{theorem}




我们可以使用 \citet{imbens2004confidence} 针对 $\tau(1,1)$ 的 confidence interval。它包含两步：第一，得到 estimated lower and upper bounds $[\hat{l},\hat{u}]$，以及 estimated standard errors $(\text{se}_l,\text{se}_u)$；第二，构造 confidence interval
$[\hat{l}-z_{1-\alpha}\text{se}_l,\hat{u}+z_{1-\alpha}\text{se}_u]$，其中 $z_{1-\alpha}$ 是 standard Normal distribution 的 $1-\alpha$ quantile。\citet{imbens2004confidence} 的 confidence interval 的有效性依赖一些 regularity conditions。在大多数 truncation by death 问题中，这些条件成立，因为 lower and upper bounds 差别明显，并且远离 extreme values $-1$ 与 $1$。这里省略技术细节。



总结一下，这是一个困难问题，因为即使 sample size 无限大，我们也不能仅根据 observed data 识别参数。我们可以为 $\tau(1,1)$ 推导 large-sample bounds，但基于 bounds 的 statistical inference 并不标准。如果没有 monotonicity，large-sample bounds 的形式还会更复杂；见 \citet{zhang2003estimation} 和 \citet[][Appendix A]{jiang2016principal}。




\subsection{An application}\label{section::ps-application}


我使用 \citet{yang2016using} 中 Acute Respiratory Distress Syndrome Network study 的数据。该研究包含 861 名 lung injury 和 acute respiratory distress syndrome 患者。患者被随机分配接受 lower tidal volumes 或 traditional tidal volumes 的 mechanical ventilation。outcome 是患者在第 28 天能否不借助辅助呼吸的二元指示变量。\cref{tb::truncationbydeath} 总结了 observed data。


\begin{table} 
\centering 
\caption{Data truncated by death，其中 * 表示死亡患者的 outcomes}\label{tb::truncationbydeath}
\begin{tabular}{cccc}
\hline 
\multicolumn{4}{c}{ Treatment $Z=1$} \\
\hline 
 & $Y=1$ & $Y=0$ &  total\\
$M=1$ &    54 & 268& 322 \\
$M=0$  & * & * & 109 \\
\hline 
\end{tabular}
~~~~~
\begin{tabular}{cccc}
\hline 
\multicolumn{4}{c}{ Control $Z=0$} \\
\hline 
 & $Y=1$ & $Y=0$ &total \\
$M=1$ &    59 & 218&  277 \\
$M=0$  & * & *& 152 \\
\hline 
\end{tabular}
\end{table}

% yang and small data
% (Z, M, Y) = (1, 1, 1): 29 + 26 = 54
% (Z, M Y) = (1, 1, 0): 258 + 10 = 268
% (Z, M, Y) = (1, 0, *): 109
% (Z, M, Y) = (0, 1, 1): 34 + 25 = 59
% (Z, M, Y) = (0, 1, 0): 211 + 7 = 218
% (Z, M, Y) = (0, 0, *): 152

我们首先得到 latent strata 的 point estimators：
\begin{eqnarray*} 
 \hat{\pi}_{(1,1)} &=& \frac{277}{277+152} = 0.646, \\
  \hat{\pi}_{(0,0)} &=&  \frac{109}{109+322} = 0.253,  \\ 
\hat{\pi}_{(1,0)} &=& 1 -0.646 -  0.253 =   0.101. 
\end{eqnarray*} 
surviving patients 的 outcome sample means 为
\begin{eqnarray*} 
\hat{E}(Y\mid Z=1,M=1) &=& \frac{54}{302} = 0.168, \\ 
\hat{E}(Y\mid Z=0, M=1) &=& \frac{59}{277} = 0.213.
\end{eqnarray*} 
$E\{Y(1)\mid M(1)=1,M(0)=1\}$ 的 bounds 估计为
$$
\left[ \frac{  ( 0.646 + 0.101 )\times  0.168-  0.101 }{  0.101 },  
\frac{ ( 0.646 + 0.101 )\times  0.168}{0.101} \right]
= [  0.037, 0.194],
$$
所以 $\tau(1,1)$ 的 bounds 为
$$
[0.037- 0.213, 0.194- 0.213] = 
[  -0.176, -0.019 ] . 
$$
基于 bootstrap 纳入 sampling uncertainty 后，使用 \citet{imbens2004confidence} 方法得到的 confidence interval 是 $[-0.267,0.039]$，它覆盖 0。




\subsection{Extensions}


\citet{zhang2003estimation} 开启了 large-sample bounds 的文献。\citet{imai2008sharp} 和 \citet{lee2009training} 是两篇后续论文。\citet{cheng2006bounds} 推导了 multiple treatment arms 下的 bounds。\citet{yang2016using} 使用 secondary outcome，\citet{yang2018using} 使用详细 survival information，来收紧 survivor average causal effect 的 bounds。

 
 
 
\section{Principal score method} 
 
 
 
 
在没有额外 assumptions 时，我们只能推出 principal strata 内 causal effects 的 bounds，但通常不能识别这些 effects。为了实现 $\tau(m_1,m_0)$ 的 nonparametric identification，我们必须施加额外 assumptions。对于这些 assumptions 的选择，目前没有共识。它们不可检验，其 plausibility 依赖具体应用。

一支研究路线基于 {\it principal score} 和 {\it principal ignorability} assumption；它与 unconfounded observational studies 中的 causal inference 平行。为简单起见，我聚焦于 strong monotonicity 的情形。



\subsection{Principal score method under strong monotonicity}

\begin{assumption}
[strong monotonicity]\label{assume::strong-monotonicity}
 $M(0) = 0$. 
\end{assumption}


类似于 ignorability assumption，现在我们假设 {\it principal ignorability} assumption。


\begin{assumption}
[principal ignorability]\label{assume::principal-ignorability-SM}
我们有
$$
E\{  Y(0) \mid M(1) = 1, X \} = E\{  Y(0) \mid M(1) = 0, X \} .
$$ 
\end{assumption}


\cref{assume::principal-ignorability-SM} 蕴含
$
E\{  Y(0) \mid M(1)  , X \} = E\{  Y(0) \mid   X \} 
$
或者等价地，
\begin{eqnarray}
\label{eq::mean-independent}
E\{  Y(0)  M(1)  \mid  X \}  =  E\{  Y(0) \mid   X \}  E\{     M(1)  \mid  X \}  ,
\end{eqnarray}
也就是说，给定 $X$ 时，$Y(0)$ 与 $M(1)$ mean independent（或 uncorrelated）。

\cref{assume::complete-randomization-withM,assume::strong-monotonicity,assume::principal-ignorability-SM} 保证 principal strata 内 causal effects 的 nonparametric identification，如下面的 \cref{thm::identification-PCEs-pscore} 所总结。


\begin{theorem}
\label{thm::identification-PCEs-pscore}
在 \cref{assume::complete-randomization-withM,assume::strong-monotonicity,assume::principal-ignorability-SM} 下，
\begin{enumerate}
[(a)]
\item 
$M(1)$ 的 conditional and marginal probabilities，
$\pi(X) =  \pr\{ M(1)=1\mid X \}  $ 和 $\pi =  \pr\{ M(1)=1\} $，分别可由
$$
\pi(X)   =  \pr(M=1\mid Z=1, X ) 
$$ 
和
$$
\pi  =   \pr(M=1\mid Z=1 ) 
$$ 
识别；
\item
principal stratification average causal effects 可由
$$
\tau(1,0) = E(Y\mid Z=1, M=1)  - E\{  \pi(X) Y  \mid Z=0 \}/\pi 
$$
和
$$
\tau(0,0) = E(Y\mid Z=1, M=0)  - E\{ (1-\pi(X)) Y  \mid Z=0 \}/(1-\pi) . 
$$
识别。
\end{enumerate}
\end{theorem}


条件概率 $\pi(X)=\pr\{M(1)=1\mid X\}$ 称为 {\it principal score}。\cref{thm::identification-PCEs-pscore} 说明，$\tau(1,0)$ 和 $\tau(0,0)$ 可以通过带有适当权重的 means difference 识别，这些权重依赖 principal score。



 \begin{myproof}{Theorem}{\cref{thm::identification-PCEs-pscore}}
 我只证明 $\tau(1,0)$ 的结果。
我们有
 $$
 E(Y\mid Z=1, M=1)  = E\{ Y(1) \mid Z=1, M(1) = 1 \} =  E\{ Y(1) \mid   M(1) = 1 \} .
 $$
此外，
 \begin{eqnarray*}
&& E  \{\pi(X) Y\mid Z=0   \}  /\pi  \\
 &=& E  \{\pi(X) Y(0)  \mid Z=0 \}  /\pi  \\
 &=&  E  \{\pi(X) Y(0)   \}  /\pi \qquad   (\text{\cref{assume::complete-randomization-withM}}) \\
 &=&  E \left[    E\{\pi(X) Y(0) \mid X  \} \right]/\pi  \qquad (\text{tower property}) \\ 
 &=& E \left[  \pi(X)  E\{ Y(0) \mid X  \} \right]/\pi  \\ 
 &=& E \left[  E\{    M(1) \mid X\} E\{ Y(0) \mid X  \} \right]/\pi \\
 &=&  E \left[  E\{    M(1)  Y(0) \mid X  \} \right]/\pi  \qquad  (\text{by \cref{eq::mean-independent}}) \\
 &=&  E\{    M(1) Y(0)  \}/\pi   \\
 &=&  E\{ Y(0) \mid M(1) = 1 \}  .
 \end{eqnarray*}
% 
%$$
%E\{ Y(0) \mid M(1) = 1 \} = 
%E\{  \pi(X) Y  \mid Z=0 \}/\pi 
%$$
%holds because 
%\begin{eqnarray*}
% E\{ Y(0) \mid M(1) = 1 \}  &=&   E\{    M(1) Y(0)  \}/\pi   \\
%&=& E \left[  E\{    M(1) \mid X\} E\{ Y(0) \mid X  \} \right]/\pi   \qquad  (\text{by \cref{eq::mean-independent}}) \\
%&=& E \left[  \pi(X)  E\{ Y(0) \mid X  \} \right]/\pi  \\
%&=& E \left[    E\{\pi(X) Y(0) \mid X  \} \right]/\pi  \\
%&=& E  \{\pi(X) Y(0)   \}  /\pi  \qquad (\text{by the tower property}) \\
%&=&  E  \{\pi(X) Y\mid Z=0   \}  /\pi. 
%\end{eqnarray*}
$\tau(0,0)$ 的结果证明留作 \cref{hw::proof-part2-identification-PCEs}。
 \end{myproof}

 
\cref{thm::identification-PCEs-pscore} 激发了下面分别针对 $\tau(1,0)$ 和 $\tau(0,0)$ 的简单 estimators：
\begin{enumerate}
\item
只使用 treated group 的数据，拟合 $M$ 对 $X$ 的 logistic regression，以得到 $\hat \pi(X_i)$；
\item
用 $\hat \pi =  \sumn Z_i M_i / \sumn Z_i$ 估计 $\pi$；
\item
得到 moment estimators：
$$
\hat \tau(1,0) =  \frac{  \sumn Z_iM_iY_i  }{ \sumn Z_iM_i }   -   \frac{   \sumn (1-Z_i) \hat \pi(X_i) Y_i  }{ \hat \pi   \sumn (1-Z_i) }   
$$
和
$$
\hat \tau(0,0) =  \frac{  \sumn Z_i(1-M_i)Y_i  }{ \sumn Z_i (1-M_i) }  - \frac{  \sumn (1-Z_i)(1-\hat \pi(X_i)\} Y_i  }{  (1- \hat \pi )\sumn (1-Z_i)  }
;
$$
\item
使用 bootstrap 近似 $\hat \tau(1,0)$ 和 $\hat \tau(0,0)$ 的 variances。
\end{enumerate}
 
 
 
 
 \subsection{An example}

下面的函数可以计算 $\hat \tau(1,0)$ 和 $\hat \tau(0,0)$ 的 point estimates。

  \begin{lstlisting}
psw = function(Z, M, Y, X) {
  ## probabilities of 10 and 00
  pi.10 = mean(M[Z==1])
  pi.00 = 1 - pi.10
  
  ## conditional probabilities of 10 and 00
  ps.10 = glm(M ~ X, family = binomial, 
              weights = Z)$fitted.values
  ps.00 = 1 - ps.10
  
  ## PCEs 10 and 00 
  tau.10 = mean(Y[Z==1&M==1]) - mean(Y[Z==0]*ps.10[Z==0])/pi.10
  tau.00 = mean(Y[Z==1&M==0]) - mean(Y[Z==0]*ps.00[Z==0])/pi.00
  c(tau.10, tau.00)
}
   \end{lstlisting}

下面的函数既可以计算 point estimators，也可以计算 bootstrap standard errors。
  \begin{lstlisting}
psw.boot = function(Z, M, Y, X, n.boot = 500){
  ## point estimates
  point.est = psw(Z, M, Y, X)
  ## bootstrap standard errors
  n = length(Z)  
  boot.est = replicate(n.boot, {
    id.boot = sample(1:n, n, replace = TRUE)
    psw(Z[id.boot], M[id.boot], Y[id.boot], X[id.boot, ])
  })
  boot.se    = apply(boot.est, 1, sd)
  ## results
  res        = rbind(point.est, boot.se)
  rownames(res) = c("est", "se")
  colnames(res) = c("tau10", "tau00")
  return(res)
}
    \end{lstlisting}
    
    
重新看 \cref{sec::cace-application} 中使用的数据。此前，我们在 IV analysis 中假设了 exclusion restriction。现在，我们去掉 exclusion restriction，但假设 principal ignorability。结果如下。


 \begin{lstlisting}
> jobsdata = read.csv("jobsdata.csv")
> getX     = lm(treat ~ sex + age + marital 
+               + nonwhite + educ + income,
+               data = jobsdata)
> X = model.matrix(getX)[, -1]
> Z = jobsdata$treat
> M = jobsdata$comply
> Y = jobsdata$job_seek
> table(Z, M)
   M
Z     0   1
  0 299   0
  1 228 372
> psw.boot(Z, M, Y, X)
        tau10       tau00
est 0.1694979 -0.09904909
se  0.1042405  0.15612983
\end{lstlisting}
        

point estimator $\hat \tau(1,0)$ 与基于 IV analysis 的估计差别不大。point estimator $\hat \tau(0,0)$ 接近 0，但 standard error 很大。在这个例子中，exclusion restriction 看起来是一个合理假设。
        
              
    
\subsection{Extensions}

 
 
 
 \citet{follmann2000effect}, \citet{hill2002differential}, \citet{jo2009use}, \citet{jo2011use} 和 \citet{stuart2015assessing} 开启了使用 principal score 识别 principal strata 内 causal effects 的文献。\citet{ding2017principal} 为这一策略提供了理论基础。他们证明了 \cref{thm::identification-PCEs-pscore} 以及 monotonicity 下的一个更一般版本；见 \cref{para::prinpscore-mono}。

 

\section{Other methods}


为了在没有 exclusion restriction assumption 的情况下估计 principal stratification average causal effects，\citet{zhang2009likelihood} 提出使用 Normal mixture models。然而，基于 Normal mixture models 的 inference 可能相当脆弱。一种策略是在某些限制下使用 additional information 来改善 inference \citep{ding2011identifiability, mealli2013using, mattei2013exploiting, jiang2016principal}。
 
 
从概念上说，principal stratification framework 适用于一般的 $M$。multi-valued $M$ 会产生许多 latent principal strata，而 continuous $M$ 会产生无穷多个 latent principal strata。在这些情形下，首先识别 principal strata 的概率就已经不平凡，更不用说识别 principal stratification average causal effects。\citet{jiang2020identification} 综述了一些有用策略。
 


\section{Homework problems}


\homeworksolutionref{26}
\paragraph{Complete the proof of \cref{thm::identification-PCEs-pscore}}\label{hw::proof-part2-identification-PCEs}

证明 \cref{thm::identification-PCEs-pscore} 中关于 $\tau(0,0)$ 的结果。


 
\paragraph{Principal score method under monotonicity}\label{para::prinpscore-mono}



本题扩展 \cref{thm::identification-PCEs-pscore}，将 \cref{assume::strong-monotonicity} 替换为 \cref{assume::monotonicity-ps}，并将 \cref{assume::principal-ignorability-SM} 替换为下面的 \cref{assume::principal-ignorability-M}。

\begin{assumption}
[principal ignorability]\label{assume::principal-ignorability-M}
我们有
$$
E\{  Y(1) \mid M(1) = 1, M(0) =0, X \} = E\{  Y(1) \mid M(1) = 1, M(0) =1, X \}
$$
以及
$$
E\{  Y(0) \mid M(1) = 1, M(0) =0, X \} = E\{  Y(0) \mid M(1) = 0, M(0) =0, X \}. 
$$
\end{assumption}


\begin{theorem}
\label{thm::identification-PCEs-pscore-general-ding-lu}
在 \cref{assume::complete-randomization-withM,assume::monotonicity-ps,assume::principal-ignorability-M} 下，
\begin{enumerate}
[(a)]
\item
 conditional and marginal principal scores 可由
\begin{eqnarray*}
\pi_{(0,0)}(X) &=& \pr(M=0\mid Z=1, X),\\
\pi_{(1,1)}(X) &=& \pr(M=1\mid Z=0, X),\\
\pi_{(1,0)}(X) &=& \pr(M=1\mid Z=1, X) - \pr(M=1\mid Z=0, X).
\end{eqnarray*}
以及
\begin{eqnarray*}
\pi_{(0,0)} &=& \pr(M=0\mid Z=1),\\
\pi_{(1,1)} &=& \pr(M=1\mid Z=0),\\
\pi_{(1,0)} &=& \pr(M=1\mid Z=1) - \pr(M=1\mid Z=0)
\end{eqnarray*}
分别识别；
\item
principal stratification average causal effects 可由
\begin{eqnarray*}
\tau(1,0) &=& E\left\{w_{1, (1,0)}(X) Y \mid Z=1, M=1\right\}  \\
&& -E\left\{w_{0, (1,0)}(X) Y \mid Z=0, M=0\right\} , \\
\tau(0,0) &=& E(Y \mid Z=1, M=0)-E\left\{w_{0, (0,0)}( X ) Y \mid Z=0, M=0\right\}, \\
\tau(1,1) &=& E\left\{w_{1,  (1,1)}( X ) Y \mid Z=1, M=1\right\}-E(Y \mid Z=0, M=1)
\end{eqnarray*}
识别，其中
\begin{eqnarray*}
w_{1, (1,0)}(X) &=&\frac{\pi_{(1,0)}(X)}{\pi_{(1,0)}(X)+\pi_{(1,1)}(X)} \Big / \frac{\pi_{(1,0)}}{\pi_{(1,0)}+\pi_{(1,1)}} , \\
w_{0, (1,0)}(X) &=& \frac{\pi_{(1,0)}(X)}{\pi_{(1,0)}(X)+\pi_{(0,0)}(X)} \Big / \frac{\pi_{(1,0)}}{\pi_{(1,0)}+\pi_{(0,0)}} ,  \\
w_{0, (0,0)}(X) &=&\frac{\pi_{(0,0)}(X)}{\pi_{(1,0)}(X)+\pi_{(0,0)}(X)} \Big / \frac{\pi_{(0,0)}}{\pi_{(1,0)}+\pi_{(0,0)}} , \\
w_{1, (1,1)}(X) &=&\frac{\pi_{(1,1)}(X)}{\pi_{(1,0)}(X)+\pi_{(1,1)}(X)} \Big / \frac{\pi_{(1,1)}}{\pi_{(1,0)}+\pi_{(1,1)}} .
\end{eqnarray*}
\end{enumerate}
\end{theorem}



Remark:
基于 \cref{thm::identification-PCEs-pscore-general-ding-lu}，我们可以构造 weighting estimators。\cref{thm::identification-PCEs-pscore-general-ding-lu} 是 \citet{ding2017principal} 中的 Proposition 2；该文还给出了 estimation 的更多细节。



\paragraph{Principal score method in observational studies}
\label{hw::principalscore-observational}

本题扩展 \cref{thm::identification-PCEs-pscore}，将 \cref{assume::complete-randomization-withM} 替换为下面的 ignorability assumption。


\begin{assumption}
\label{assume::ignorabilitywithM}
$
Z\ind \{ M(1), M(0), Y(1) , Y(0) \} \mid X.
$
\end{assumption}

回忆 propensity score 的定义 $e(X)=\pr(Z=1\mid X)$。
我们有下面的 identification result。


\begin{theorem}
\label{thm::identification-PCEs-pscore-obs}
在 \cref{assume::ignorabilitywithM,assume::strong-monotonicity,assume::principal-ignorability-SM} 下，
\begin{enumerate}
[(a)]
\item 
$M(1)$ 的 conditional and marginal probabilities，
$\pi(X) =  \pr\{ M(1)=1\mid X \}  $ 和 $\pi =  \pr\{ M(1)=1\} $，分别可由
$$
\pi(X)   =  \pr(M=1\mid Z=1, X ) 
$$ 
以及
$$
\pi  =  E\{\pr(M=1\mid Z=1, X ) \}
$$ 
识别；
\item
principal stratification average causal effects 可由
$$
\tau(1,0) = E\left\{   \frac{M}{\pi}   \frac{Z}{e(X)} Y \right\}
- E\left\{   \frac{ \pi(X)  }{\pi}   \frac{1-Z}{1-e(X)} Y    \right\} 
$$
和
$$
\tau(0,0) = E\left\{   \frac{1-M}{1-\pi}   \frac{Z}{e(X)} Y \right\}
- E\left\{   \frac{ 1- \pi(X)  }{1 - \pi}   \frac{1-Z}{1-e(X)} Y    \right\} .
$$
识别。
\end{enumerate}
\end{theorem}


证明 \cref{thm::identification-PCEs-pscore-obs}。


 


\paragraph{General principal score method}
\label{hw::general-principal-score-jiangetal}

在 \cref{assume::ignorabilitywithM,assume::monotonicity-ps,assume::principal-ignorability-M} 下，扩展 \cref{thm::identification-PCEs-pscore-general-ding-lu,thm::identification-PCEs-pscore-obs}。


Remark: 见 \citet{jiang2020multiply}.


  
 
